\textbf{Visual Language Navigation.} 
Visual-Language Navigation~\cite{anderson2018vision} requires an agent to navigate in unknown photo-realistic environments based on its visual perceptions following natural language instructions. Navigation agents should be capable of capturing spatiotemporal information such as the relative motion of objects from navigation histories, especially when the agent navigates with short step size in continuous spaces. To verify the effectiveness of such ability of our model, we conduct our experiments on the VLN-CE benchmark~\cite{krantz2020beyond}, demanding the agent to function in a continuous environment.

We conduct our experiments using the method proposed in~\cite{an20221st}(CWP-HEP). The history-enhanced planner is a customized variant of HAMT~\cite{chen2021history} which uses a concatenation of depth embedding and RGB embedding as the input embedding. Note that we don't use the tryout controller here since VLN-CE setting allows sliding. This is a strong baseline that already outperforms the previous state-of-the-art method CWP-VLNBERT~\cite{hong2022bridging}. In each decision loop, we collect the latest 16-frame observations to form panoramic navigation videos and then encode the video using ViT-L in \model. The video embedding is concatenated with the RGB embedding and depth embedding as the final image embedding. For evaluation, we refer to~\cite{an20221st} for detailed metrics. \model~could improved our baseline from 50.2\% to 52.9\% in Success Rate(SR) (Table \ref{tab:vlnce}).